% =====================================================================
%  BÁO CÁO KIỂM THỬ - USECASE: CHỈNH SỬA YÊU CẦU NHẬP HÀNG
% ---------------------------------------------------------------------
%  HƯỚNG DẪN BIÊN DỊCH (QUAN TRỌNG):
%   - Biên dịch bằng XeLaTeX (KHÔNG dùng pdfLaTeX vì có tiếng Việt Unicode).
%   - Local (MiKTeX/TeX Live):   xelatex BaoCaoKiemThu_YeuCauNhapHang.tex
%                                 (chạy 2 lần để cập nhật Mục lục)
%   - Overleaf:  Menu -> Compiler -> chọn "XeLaTeX".
%  Font dùng "TeX Gyre Termes" (giống Times New Roman, có sẵn glyph
%  tiếng Việt, đi kèm mọi bản TeX) nên không cần cài thêm.
% =====================================================================
\documentclass[12pt,a4paper]{article}

% ----- Font & Unicode (XeLaTeX) -----
\usepackage{fontspec}
\setmainfont{TeX Gyre Termes}      % Times-like, hỗ trợ tiếng Việt
\setsansfont{TeX Gyre Heros}
% \setmonofont{...}  % để mặc định Latin Modern Mono cho phần code (ASCII)

% ----- Trang & bố cục -----
\usepackage{geometry}
\geometry{a4paper, margin=2.5cm}
\usepackage{amsmath,amssymb}
\usepackage{graphicx}
\usepackage{float}
\usepackage[table]{xcolor}
\usepackage{array}
\usepackage{tabularx}
\usepackage{longtable}
\usepackage{pdflscape}
\usepackage{enumitem}
\usepackage{listings}
\usepackage{fancyhdr}
\usepackage{tikz}
\usetikzlibrary{shapes.geometric, arrows.meta, positioning}
\usepackage[hidelinks]{hyperref}
\hypersetup{unicode=true, colorlinks=true, linkcolor=blue!60!black,
            urlcolor=blue!60!black, citecolor=blue!60!black}

% ----- Tên nhãn tiếng Việt (thay cho babel để giảm phụ thuộc) -----
\renewcommand{\contentsname}{Mục lục}
\renewcommand{\figurename}{Hình}
\renewcommand{\tablename}{Bảng}
\renewcommand{\refname}{Tài liệu tham khảo}
% Ngắt dòng tiếng Việt tại khoảng trắng (mặc định đã ổn, thêm cho chắc)
\XeTeXlinebreaklocale "vi"
\XeTeXlinebreakskip = 0pt plus 0.1pt

% ----- Kiểu cột bảng -----
\newcolumntype{Y}{>{\raggedright\arraybackslash}X}
\newcolumntype{C}{>{\centering\arraybackslash}X}
\renewcommand{\arraystretch}{1.25}

% ----- Header / Footer -----
\pagestyle{fancy}
\fancyhf{}
\lhead{\small Kiểm thử: Chỉnh sửa yêu cầu nhập hàng}
\rhead{\small Nhóm 4}
\cfoot{\thepage}
\renewcommand{\headrulewidth}{0.4pt}

% ----- Kiểu code Java -----
\definecolor{codekw}{rgb}{0.13,0.13,0.7}
\definecolor{codecm}{rgb}{0.0,0.45,0.0}
\definecolor{codestr}{rgb}{0.6,0.1,0.1}
\definecolor{codebg}{rgb}{0.97,0.97,0.97}
\lstset{
  language=Java,
  basicstyle=\ttfamily\footnotesize,
  keywordstyle=\color{codekw}\bfseries,
  commentstyle=\color{codecm}\itshape,
  stringstyle=\color{codestr},
  numbers=left, numberstyle=\tiny\color{gray}, numbersep=8pt,
  backgroundcolor=\color{codebg},
  frame=single, rulecolor=\color{gray!50},
  breaklines=true, showstringspaces=false, tabsize=2,
  captionpos=b, xleftmargin=14pt, framexleftmargin=14pt
}

% ----- Thông tin sinh viên (CHỈNH SỬA TẠI ĐÂY NẾU CẦN) -----
\newcommand{\svName}{Phú Hưng}          % TODO: kiểm tra lại họ tên đầy đủ
\newcommand{\svMSSV}{20235738}
\newcommand{\svGVHD}{Trịnh Tuấn Đạt}
\newcommand{\svHP}{IT4549}
\newcommand{\svNhom}{Nhóm 4}

\begin{document}

% =====================================================================
%  TRANG BÌA
% =====================================================================
\begin{titlepage}
\centering
\thispagestyle{empty}

{\large \textbf{ĐẠI HỌC BÁCH KHOA HÀ NỘI}}\\[2pt]
{\large \textbf{TRƯỜNG CÔNG NGHỆ THÔNG TIN VÀ TRUYỀN THÔNG}}\\[6pt]
\rule{\textwidth}{0.8pt}

\vspace{1.2cm}
% Nếu có logo, bỏ ghi chú dòng dưới và đặt file logo.png cùng thư mục:
% \includegraphics[width=3.5cm]{logo.png}\\[1cm]

\vspace{0.8cm}
{\Huge \textbf{BÁO CÁO KIỂM THỬ}}\\[10pt]
{\LARGE \textbf{HỆ THỐNG ĐẶT HÀNG NHẬP KHẨU}}\\[18pt]
{\Large Usecase: \textbf{Chỉnh sửa yêu cầu nhập hàng}}\\[6pt]
{\large (Thiết kế \& thực thi kiểm thử: Hộp đen -- Hộp trắng C1 -- Use Case)}

\vspace{1.6cm}
\begin{flushleft}\large
\hspace{1.5cm}
\begin{tabular}{@{}ll@{}}
\textbf{Giảng viên hướng dẫn:} & \svGVHD \\[8pt]
\textbf{Sinh viên thực hiện:}  & \svName \\[8pt]
\textbf{MSSV:}                 & \svMSSV \\[8pt]
\textbf{Mã số lớp học phần:}   & \svHP \\[8pt]
\textbf{Nhóm:}                 & \svNhom \\
\end{tabular}
\end{flushleft}

\vfill
{\large Hà Nội, 2026}
\end{titlepage}

% =====================================================================
%  MỤC LỤC
% =====================================================================
\tableofcontents
\thispagestyle{empty}
\newpage
\setcounter{page}{1}

% =====================================================================
\section{Mô tả phương thức/lớp cần kiểm thử}

\subsection{Thông tin lớp và phương thức kiểm thử đơn vị}
\begin{itemize}[leftmargin=*]
  \item \textbf{Tên lớp kiểm thử đơn vị:} \texttt{com.app.modules.sales.request.editrequest.service.RequestService}
  \item \textbf{Tên phương thức kiểm thử đơn vị:} \texttt{updateRequest(UpdateRequestDTO dto)}
  \item \textbf{Chức năng nghiệp vụ:} Phương thức chịu trách nhiệm cập nhật các thông tin thay đổi của một yêu cầu nhập hàng từ giao diện UI gửi xuống CSDL. Quy trình xử lý gồm:
  \begin{enumerate}[leftmargin=*]
    \item Kiểm tra sự tồn tại của yêu cầu nhập hàng bằng mã yêu cầu (\texttt{dto.getCode()}).
    \item Kiểm tra trạng thái: chỉ cho phép chỉnh sửa nếu yêu cầu đang ở trạng thái \texttt{pending}. Nếu ở trạng thái khác (\texttt{processing} hoặc \texttt{completed}), phương thức ném ngoại lệ \texttt{IllegalStateException}.
    \item Đồng bộ hóa danh sách mặt hàng (upsert): duyệt danh sách mặt hàng mới từ DTO, gọi repository để thêm mới hoặc cập nhật số lượng/ngày nhận của từng mặt hàng.
    \item Xóa bỏ các mặt hàng cũ trong yêu cầu mà không còn xuất hiện trong danh sách mới của DTO.
    \item Tải lại thông tin yêu cầu mới nhất từ CSDL và trả về đối tượng \texttt{RequestResponse}.
  \end{enumerate}
\end{itemize}

\subsection{Chữ ký phương thức (Method Signature)}
\begin{lstlisting}
public RequestResponse updateRequest(UpdateRequestDTO dto)
\end{lstlisting}

\begin{itemize}[leftmargin=*]
  \item \textbf{Tham số đầu vào:} \texttt{UpdateRequestDTO dto} -- chứa mã yêu cầu cần chỉnh sửa (\texttt{code}) và danh sách mặt hàng (\texttt{List<RequestItem> items}) biểu diễn trạng thái mới mà người dùng mong muốn.
  \item \textbf{Giá trị trả về:} \texttt{RequestResponse} -- DTO chứa thông tin chi tiết yêu cầu và danh sách mặt hàng đã được lưu thành công trong CSDL.
  \item \textbf{Ngoại lệ:}
  \begin{itemize}
    \item \texttt{NoSuchElementException}: khi mã yêu cầu trong DTO không tìm thấy trong CSDL.
    \item \texttt{IllegalStateException}: khi yêu cầu không ở trạng thái \texttt{pending}.
  \end{itemize}
\end{itemize}

% =====================================================================
\section{Thiết kế test case Hộp đen (Black-box Testing)}

Kỹ thuật kiểm thử hộp đen được áp dụng là \textbf{Phân vùng tương đương (Equivalence Partitioning)} kết hợp \textbf{Phân tích giá trị biên (Boundary Value Analysis)} đối với các tham số đầu vào và trạng thái nghiệp vụ.

\subsection{Phân hoạch các lớp tương đương}
\begin{table}[H]
\centering
\footnotesize
\begin{tabularx}{\textwidth}{|p{3.4cm}|Y|Y|}
\hline
\rowcolor{gray!20}
\textbf{Tham số / Trạng thái} & \textbf{Lớp tương đương hợp lệ} & \textbf{Lớp tương đương không hợp lệ}\\ \hline
Mã yêu cầu (Request Code) & Mã yêu cầu tồn tại trong hệ thống (H1) & Mã yêu cầu không tồn tại (KH1)\\ \hline
Trạng thái yêu cầu (Status) & Trạng thái là \texttt{pending} (H2) & Trạng thái \texttt{processing} (KH2); Trạng thái \texttt{completed} (KH3)\\ \hline
Số lượng mặt hàng trong DTO & Số mặt hàng $\geq 1$ (H3); Số mặt hàng $= 0$ -- xóa sạch (H4) & (không có) \\ \hline
\end{tabularx}
\end{table}

\subsection{Danh sách các Test Case Hộp đen}
\noindent{\footnotesize
\begin{longtable}{|p{1.5cm}|p{4.0cm}|p{4.5cm}|p{4.6cm}|}
\hline
\rowcolor{gray!20}
\textbf{Mã TC} & \textbf{Mục tiêu kiểm thử} & \textbf{Đầu vào / Trạng thái hệ thống} & \textbf{Kết quả mong đợi}\\ \hline
\endfirsthead
\hline
\rowcolor{gray!20}
\textbf{Mã TC} & \textbf{Mục tiêu kiểm thử} & \textbf{Đầu vào / Trạng thái hệ thống} & \textbf{Kết quả mong đợi}\\ \hline
\endhead
\textbf{TC\_BB\_01} & Cập nhật thành công yêu cầu \texttt{pending} với danh sách hợp lệ (thêm mới, giữ lại/cập nhật, xóa bớt). & RQ001 (\texttt{pending}, 2 mặt hàng cũ I01, I02); DTO: RQ001 với I01 (đổi SL/ngày) và I03 (mới). & Lưu thành công; trả về \texttt{RequestResponse} chứa 2 mặt hàng I01, I03; I02 bị xóa khỏi hệ thống.\\ \hline
\textbf{TC\_BB\_02} & Báo lỗi khi mã yêu cầu trong DTO không tồn tại. & Không tồn tại RQ999; DTO: RQ999, danh sách rỗng. & Ném \texttt{NoSuchElementException} với thông báo ``Không tìm thấy yêu cầu RQ999''.\\ \hline
\textbf{TC\_BB\_03} & Ngăn chỉnh sửa khi yêu cầu đang \texttt{processing}. & RQ002 (\texttt{processing}); DTO: RQ002. & Ném \texttt{IllegalStateException}: ``Yêu cầu RQ002 đang ở trạng thái processing, không thể chỉnh sửa.''\\ \hline
\textbf{TC\_BB\_04} & Ngăn chỉnh sửa khi yêu cầu đã \texttt{completed}. & RQ003 (\texttt{completed}); DTO: RQ003. & Ném \texttt{IllegalStateException}: ``Yêu cầu RQ003 đang ở trạng thái completed, không thể chỉnh sửa.''\\ \hline
\textbf{TC\_BB\_05} & Cho phép xóa sạch toàn bộ mặt hàng (giá trị biên: danh sách rỗng). & RQ004 (\texttt{pending}, 1 mặt hàng cũ I01); DTO: RQ004, danh sách rỗng. & Cập nhật thành công; trả về \texttt{RequestResponse} có số mặt hàng $= 0$; I01 bị xóa khỏi DB.\\ \hline
\end{longtable}
}

% =====================================================================
\section{Thiết kế test case Hộp trắng (White-box -- Độ đo C1)}

Kỹ thuật kiểm thử hộp trắng được áp dụng nhằm đảm bảo \textbf{Độ đo bao phủ nhánh (C1 -- Branch Coverage)} đạt mức tối đa 100\%.

\subsection{Mã nguồn và Đồ thị dòng kiểm soát (CFG)}
Phương thức \texttt{updateRequest} (đã lược bỏ chú thích) được kiểm thử:

\begin{lstlisting}
public RequestResponse updateRequest(UpdateRequestDTO dto) {
    Request request = findRequestOrThrow(dto.getCode());          // Dinh 1
    if (!"pending".equals(request.getStatus())) {                 // Quyet dinh 1
        throw new IllegalStateException("...");                   // Dinh 3
    }
    Set<String> keep = new HashSet<>();                           // Dinh 4
    int pos = 0;
    for (RequestItem item : dto.getItems()) {                     // Quyet dinh 2
        repository.insertItem(dto.getCode(), item, pos++);        // Dinh 6
        repository.updateItemDeliveryDate(dto.getCode(),
                item.getCode(), item.getDeliveryDate());
        keep.add(item.getCode());
    }
    for (RequestItem old : request.getItems()) {                  // Quyet dinh 3
        if (!keep.contains(old.getCode())) {                      // Quyet dinh 4
            repository.deleteItem(dto.getCode(), old.getCode());  // Dinh 9
        }
    }
    return new RequestResponse(findRequestOrThrow(dto.getCode()));// Dinh 11
}
\end{lstlisting}

\begin{figure}[H]
\centering
\begin{tikzpicture}[
  node distance = 9mm and 16mm,
  start/.style   = {ellipse, draw, fill=gray!12, align=center, font=\footnotesize, minimum width=24mm},
  block/.style   = {rectangle, draw, fill=blue!6, align=center, font=\scriptsize, text width=42mm, inner sep=4pt},
  dec/.style     = {diamond, draw, fill=orange!14, aspect=2.4, align=center, font=\scriptsize, inner sep=0pt},
  fl/.style      = {-{Latex[length=2mm]}, semithick}
]
\node[start] (st) {Bắt đầu};
\node[block, below=of st]  (n1) {Đỉnh 1: tìm yêu cầu\\ \texttt{findRequestOrThrow}};
\node[dec,   below=of n1]  (n2) {Đỉnh 2:\\ status $\neq$ pending?};
\node[block, right=22mm of n2] (n3) {Đỉnh 3: ném\\ \texttt{IllegalStateException}};
\node[start, below=of n3]  (e1) {Kết thúc (ngoại lệ)};
\node[block, below=12mm of n2] (n4) {Đỉnh 4: khởi tạo\\ \texttt{keep}, \texttt{pos=0}};
\node[dec,   below=of n4]  (n5) {Đỉnh 5:\\ còn item DTO?};
\node[block, right=20mm of n5] (n6) {Đỉnh 6: \texttt{insertItem} +\\ cập nhật ngày + \texttt{keep.add}};
\node[dec,   below=14mm of n5] (n7) {Đỉnh 7:\\ còn item cũ?};
\node[dec,   below=13mm of n7] (n8) {Đỉnh 8:\\ \texttt{!keep.contains}?};
\node[block, right=20mm of n8] (n9) {Đỉnh 9:\\ \texttt{deleteItem}};
\node[block, below=16mm of n8] (n11){Đỉnh 11: trả về\\ \texttt{RequestResponse}};
\node[start, below=of n11] (fn) {Kết thúc (thành công)};

\draw[fl] (st) -- (n1);
\draw[fl] (n1) -- (n2);
\draw[fl] (n2) -- node[above,font=\scriptsize]{Đúng (1.A)} (n3);
\draw[fl] (n3) -- (e1);
\draw[fl] (n2) -- node[right,font=\scriptsize]{Sai (1.B)} (n4);
\draw[fl] (n4) -- (n5);
\draw[fl] (n5) -- node[above,font=\scriptsize]{Có (2.A)} (n6);
\draw[fl] (n6) to[bend left=55] (n5);
\draw[fl] (n5) -- node[right,font=\scriptsize]{Hết (2.B)} (n7);
\draw[fl] (n7) -- node[right,font=\scriptsize]{Có (3.A)} (n8);
\draw[fl] (n8) -- node[above,font=\scriptsize]{Đúng (4.A)} (n9);
\draw[fl] (n9) to[bend left=55] (n7);
\draw[fl] (n8) to[bend right=70] node[left,font=\scriptsize]{Sai (4.B)} (n7);
\draw[fl] (n7) -- node[right,font=\scriptsize]{Hết (3.B)} (n11);
\draw[fl] (n11) -- (fn);
\end{tikzpicture}
\caption{Đồ thị dòng kiểm soát (CFG) của phương thức \texttt{updateRequest}.}
\end{figure}

\subsection{Các quyết định và nhánh điều hướng}
Có 4 điểm quyết định tương đương 8 nhánh điều hướng cần bao phủ:
\begin{enumerate}[leftmargin=*]
  \item \textbf{Quyết định 1} (\texttt{!"pending".equals(status)}): Nhánh \textbf{1.A} (True -- trạng thái khác \texttt{pending}); Nhánh \textbf{1.B} (False -- bằng \texttt{pending}).
  \item \textbf{Quyết định 2} (duyệt vòng lặp DTO): Nhánh \textbf{2.A} (còn phần tử); Nhánh \textbf{2.B} (hết/không có phần tử).
  \item \textbf{Quyết định 3} (duyệt danh sách mặt hàng cũ): Nhánh \textbf{3.A} (còn phần tử); Nhánh \textbf{3.B} (hết/không có phần tử).
  \item \textbf{Quyết định 4} (\texttt{!keep.contains(old.getCode())}): Nhánh \textbf{4.A} (True -- mặt hàng cũ bị xóa); Nhánh \textbf{4.B} (False -- mặt hàng cũ được giữ lại).
\end{enumerate}

\subsection{Thiết kế Test Case và Bảng phủ nhánh (C1 Coverage Matrix)}
Để phủ 100\% các nhánh, thiết kế 4 test case hộp trắng:
\begin{itemize}[leftmargin=*]
  \item \textbf{TC\_WB\_01:} phủ nhánh ngoại lệ trạng thái (1.A).
  \item \textbf{TC\_WB\_02:} phủ các nhánh lặp trống (1.B, 2.B, 3.B).
  \item \textbf{TC\_WB\_03:} phủ toàn bộ nhánh lặp và điều kiện bên trong (1.B, 2.A, 2.B, 3.A, 3.B, 4.A, 4.B).
  \item \textbf{TC\_WB\_04:} trường hợp đặc biệt không tìm thấy yêu cầu (ném \texttt{NoSuchElementException}).
\end{itemize}

\begin{table}[H]
\centering
\footnotesize
\begin{tabularx}{\textwidth}{|Y|c|c|c|c|}
\hline
\rowcolor{gray!20}
\textbf{Nhánh điều hướng} & \textbf{TC\_WB\_01} & \textbf{TC\_WB\_02} & \textbf{TC\_WB\_03} & \textbf{TC\_WB\_04}\\ \hline
Nhánh 1.A (trạng thái khác pending) & \textbf{x} & & & \\ \hline
Nhánh 1.B (trạng thái pending) & & \textbf{x} & \textbf{x} & \\ \hline
Nhánh 2.A (DTO có phần tử) & & & \textbf{x} & \\ \hline
Nhánh 2.B (DTO không còn phần tử) & & \textbf{x} & \textbf{x} & \\ \hline
Nhánh 3.A (hàng cũ có phần tử) & & & \textbf{x} & \\ \hline
Nhánh 3.B (hàng cũ không còn phần tử) & & \textbf{x} & \textbf{x} & \\ \hline
Nhánh 4.A (hàng cũ bị xóa) & & & \textbf{x} & \\ \hline
Nhánh 4.B (hàng cũ được giữ lại) & & & \textbf{x} & \\ \hline
Ngoại lệ NoSuchElement & & & & \textbf{x}\\ \hline
\end{tabularx}
\end{table}

% =====================================================================
\section{Cài đặt chương trình kiểm thử tự động}

Các test case trên được lập trình tự động bằng Java, sử dụng \textbf{JUnit 5 (Jupiter)}.
\begin{itemize}[leftmargin=*]
  \item \textbf{Tên đầy đủ của Class kiểm thử (Full Name):}\\ \texttt{com.app.modules.sales.request.editrequest.service.RequestServiceTest}
  \item \textbf{Mã nguồn lớp kiểm thử:} \texttt{src/test/java/com/app/modules/sales/request/editrequest/service/RequestServiceTest.java}
  \item \textbf{Phương pháp Stubbing:} dùng lớp nội bộ tĩnh \texttt{FakeRequestRepository} kế thừa \texttt{RequestRepository} để ghi nhận tương tác và cung cấp dữ liệu giả lập trong bộ nhớ, không phụ thuộc CSDL PostgreSQL, tránh phản tác dụng và đảm bảo tốc độ chạy nhanh.
  \item \textbf{Kết quả thực thi (đã chạy thực tế bằng \texttt{mvn test}):}
\end{itemize}
\begin{lstlisting}[language=bash, numbers=none]
Tests run: 9, Failures: 0, Errors: 0, Skipped: 0  -- 9/9 PASS
\end{lstlisting}

% =====================================================================
\section{Kiểm thử Use Case (Use Case Testing)}

Khác với kiểm thử đơn vị (kiểm thử riêng phương thức \texttt{updateRequest}), phần này kiểm thử ở \textbf{mức hệ thống (system/acceptance level)} cho toàn bộ use case ``Chỉnh sửa yêu cầu nhập hàng'' trên giao diện. Quy trình áp dụng template gồm 4 bước: \emph{(1) Đặc tả Use Case $\to$ (2) Xác định kịch bản (Scenarios) $\to$ (3) Thiết kế Test Case từ kịch bản $\to$ (4) Thực thi và ghi nhận kết quả.}

\subsection{Đặc tả Use Case (Use Case Specification)}
\begin{table}[H]
\centering
\small
\begin{tabularx}{\textwidth}{|>{\bfseries}p{4.2cm}|Y|}
\hline
\rowcolor{gray!20}\normalfont\textbf{Mục} & \textbf{Nội dung}\\ \hline
Mã Use Case & UC-SALES-EDIT-01\\ \hline
Tên Use Case & Chỉnh sửa yêu cầu nhập hàng\\ \hline
Tác nhân chính (Actor) & Nhân viên kinh doanh / Bộ phận bán hàng\\ \hline
Mô tả ngắn & Cho phép nhân viên kinh doanh cập nhật danh sách mặt hàng (số lượng, ngày nhận, thêm/xóa mặt hàng) của một yêu cầu nhập hàng đang ở trạng thái chờ duyệt và lưu thay đổi xuống CSDL.\\ \hline
Điều kiện kích hoạt (Trigger) & Người dùng chọn một yêu cầu và mở màn hình ``Chỉnh sửa yêu cầu''.\\ \hline
Tiền điều kiện & Đã đăng nhập với quyền bộ phận bán hàng; yêu cầu cần sửa đã tồn tại trong hệ thống.\\ \hline
Hậu điều kiện thành công & Bảng \texttt{RequestDetail} được cập nhật đúng trong CSDL; cờ thay đổi (dirty) được đặt lại $=$ false.\\ \hline
Hậu điều kiện thất bại & Không có dữ liệu nào bị thay đổi trong CSDL.\\ \hline
\end{tabularx}
\end{table}

\noindent\textbf{Luồng sự kiện chính (Main Flow -- MF):}
\begin{enumerate}[leftmargin=*]
  \item Hệ thống tải chi tiết yêu cầu và hiển thị thông tin chung + bảng mặt hàng.
  \item Yêu cầu ở trạng thái \texttt{pending} $\Rightarrow$ hệ thống bật nút ``Thêm mặt hàng'', ``Lưu thay đổi'' và cho phép sửa ô Số lượng, Ngày nhận, hiển thị nút Xóa.
  \item Người dùng sửa Số lượng của một mặt hàng; hệ thống đánh dấu ``có thay đổi'' (dirty = true).
  \item Người dùng sửa Ngày nhận của một mặt hàng (dirty = true).
  \item Người dùng bấm ``Thêm mặt hàng'' $\to$ chọn sản phẩm $\to$ nhập số lượng ($> 0$) và ngày nhận hợp lệ $\to$ mặt hàng mới được thêm vào bảng (dirty = true).
  \item Người dùng bấm nút Xóa ở một mặt hàng $\to$ xác nhận ``OK'' $\to$ mặt hàng bị gỡ khỏi bảng (dirty = true).
  \item Người dùng bấm ``Lưu thay đổi''.
  \item Hệ thống gọi service đồng bộ dữ liệu xuống CSDL, hiển thị ``Đã lưu thay đổi yêu cầu [mã]'', đặt dirty = false. \textbf{Use case kết thúc thành công.}
\end{enumerate}

\noindent\textbf{Luồng thay thế (Alternative Flows -- AF):}
\begin{itemize}[leftmargin=*]
  \item \textbf{AF1 -- Hủy xác nhận xóa:} ở bước 6, người dùng chọn ``Cancel'' $\Rightarrow$ mặt hàng KHÔNG bị xóa, quay lại bước 6.
  \item \textbf{AF2 -- Quay lại \& chấp nhận hủy thay đổi:} khi đang có thay đổi chưa lưu, bấm ``Quay lại'' $\Rightarrow$ hộp thoại ``Thoát mà không lưu'' $\Rightarrow$ chọn ``OK'' $\Rightarrow$ rời màn hình, mọi thay đổi bị hủy.
  \item \textbf{AF3 -- Quay lại \& ở lại:} như AF2 nhưng chọn ``Cancel'' $\Rightarrow$ đóng hộp thoại, ở lại màn hình, giữ nguyên thay đổi.
  \item \textbf{AF4 -- Quay lại khi không có thay đổi:} dirty = false $\Rightarrow$ không hỏi, rời màn hình ngay.
\end{itemize}

\noindent\textbf{Luồng ngoại lệ (Exception Flows -- EF):}
\begin{itemize}[leftmargin=*]
  \item \textbf{EF1 -- Yêu cầu không ở trạng thái pending:} nếu trạng thái là \texttt{processing}/\texttt{completed} (editable = false) $\Rightarrow$ thông báo ``Không thể chỉnh sửa'', ẩn nút ``Thêm mặt hàng'', vô hiệu hóa nút ``Lưu'', các ô Số lượng/Ngày nhận chỉ đọc, ẩn nút Xóa.
  \item \textbf{EF2 -- Cố lưu khi không được phép sửa:} nếu hành động Lưu bị kích hoạt khi editable = false $\Rightarrow$ cảnh báo ``Yêu cầu [mã] không còn ở trạng thái PENDING, không thể lưu thay đổi.'', không ghi CSDL.
  \item \textbf{EF3 -- Số lượng không hợp lệ khi sửa trực tiếp trên bảng:} nhập số âm $\Rightarrow$ tự kẹp về 0; nhập ký tự không phải số $\Rightarrow$ khôi phục giá trị cũ. Không bao giờ ghi nhận giá trị không hợp lệ.
  \item \textbf{EF4 -- Dữ liệu không hợp lệ ở hộp thoại ``Thêm mặt hàng'':} số lượng $\leq 0$ $\Rightarrow$ ``Số lượng phải là số nguyên dương.''; chưa chọn ngày $\Rightarrow$ ``Vui lòng chọn ngày nhận mong muốn.''; không thêm được mặt hàng.
  \item \textbf{EF5 -- Không tải được yêu cầu:} nếu service ném lỗi (vd mã không tồn tại) $\Rightarrow$ cảnh báo ``Không thể tải yêu cầu'' kèm ``Không tìm thấy yêu cầu [mã]''.
\end{itemize}

\subsection{Xác định các kịch bản kiểm thử (Test Scenarios)}
Mỗi kịch bản (scenario) là một đường đi (path) xuyên qua use case, kết hợp luồng chính với một luồng thay thế/ngoại lệ. Từ đặc tả trên, xác định được \textbf{13 kịch bản}:

\noindent{\small
\begin{longtable}{|p{1.6cm}|p{7.0cm}|p{3.6cm}|p{1.9cm}|}
\hline
\rowcolor{gray!20}
\textbf{Mã KB} & \textbf{Mô tả đường đi} & \textbf{Luồng kết hợp} & \textbf{Loại}\\ \hline
\endfirsthead
\hline\rowcolor{gray!20}
\textbf{Mã KB} & \textbf{Mô tả đường đi} & \textbf{Luồng kết hợp} & \textbf{Loại}\\ \hline
\endhead
SC-01 & Sửa SL/ngày, thêm 1 mặt hàng, xóa 1 mặt hàng rồi lưu thành công & MF (1$\to$8) & Chính\\ \hline
SC-02 & Bấm xóa mặt hàng nhưng hủy xác nhận & MF(1-5) + AF1 & Thay thế\\ \hline
SC-03 & Có thay đổi, bấm Quay lại, đồng ý hủy thay đổi & MF(1-3) + AF2 & Thay thế\\ \hline
SC-04 & Có thay đổi, bấm Quay lại, chọn ở lại & MF(1-3) + AF3 & Thay thế\\ \hline
SC-05 & Không thay đổi gì, bấm Quay lại & AF4 & Thay thế\\ \hline
SC-06 & Mở yêu cầu trạng thái \texttt{processing} (chỉ xem) & EF1 & Ngoại lệ\\ \hline
SC-07 & Mở yêu cầu trạng thái \texttt{completed} (chỉ xem) & EF1 & Ngoại lệ\\ \hline
SC-08 & Cố lưu khi yêu cầu không còn pending & EF2 & Ngoại lệ\\ \hline
SC-09 & Sửa số lượng $=$ số âm trên bảng & EF3 & Ngoại lệ\\ \hline
SC-10 & Sửa số lượng $=$ chữ trên bảng & EF3 & Ngoại lệ\\ \hline
SC-11 & Thêm mặt hàng với số lượng $\leq 0$ & EF4 & Ngoại lệ\\ \hline
SC-12 & Thêm mặt hàng nhưng chưa chọn ngày nhận & EF4 & Ngoại lệ\\ \hline
SC-13 & Mở yêu cầu có mã không tồn tại & EF5 & Ngoại lệ\\ \hline
\end{longtable}
}

\noindent$\Rightarrow$ \textbf{Số lượng test case cần thiết kế: 13} (mỗi kịch bản tối thiểu 1 test case; hai luồng ngoại lệ EF3 và EF4 được tách thành 2 test case theo từng loại dữ liệu lỗi để kiểm chứng riêng từng nhánh xử lý).

\subsection{Thiết kế Test Case cho Use Case}
\textit{Ghi chú: trong hệ thống thật, mã yêu cầu là ID dạng số trong CSDL nên các test case dưới đây dùng mã mẫu dạng số (ví dụ 1001).}

\begin{landscape}
\thispagestyle{empty}
\noindent{\footnotesize
\begin{longtable}{|p{1.4cm}|p{1.2cm}|p{3.4cm}|p{7.2cm}|p{5.4cm}|p{2.6cm}|p{1.4cm}|}
\hline
\rowcolor{gray!20}
\textbf{Mã TC} & \textbf{Kịch bản} & \textbf{Tiền điều kiện} & \textbf{Các bước thực hiện (kèm dữ liệu)} & \textbf{Kết quả mong đợi} & \textbf{Kết quả thực tế} & \textbf{Trạng thái}\\ \hline
\endfirsthead
\hline\rowcolor{gray!20}
\textbf{Mã TC} & \textbf{Kịch bản} & \textbf{Tiền điều kiện} & \textbf{Các bước thực hiện (kèm dữ liệu)} & \textbf{Kết quả mong đợi} & \textbf{Kết quả thực tế} & \textbf{Trạng thái}\\ \hline
\endhead
\textbf{UC-TC-01} & SC-01 & Yêu cầu 1001 ở trạng thái \texttt{pending}, có $\geq 2$ mặt hàng & (1) Mở màn chỉnh sửa 1001. (2) Sửa SL mặt hàng đầu $\to$ 120, Enter. (3) Đổi ngày nhận. (4) ``Thêm mặt hàng'', chọn 1 SP, nhập SL=30 + ngày hợp lệ, Xác nhận. (5) Xóa 1 mặt hàng cũ, chọn OK. (6) ``Lưu thay đổi''. & Thông báo ``Đã lưu thay đổi yêu cầu 1001''; bảng phản ánh đúng (mặt hàng sửa có SL/ngày mới, có mặt hàng mới, mặt hàng bị xóa biến mất); dirty=false; CSDL đúng. & Đúng như mong đợi & \textbf{PASS}\\ \hline
\textbf{UC-TC-02} & SC-02 & Yêu cầu 1001 \texttt{pending}, có $\geq 1$ mặt hàng & (1) Bấm nút Xóa ở 1 mặt hàng. (2) Hộp thoại ``Xác nhận xóa mặt hàng'' $\to$ chọn ``Cancel''. & Hộp thoại đóng; mặt hàng vẫn còn trong bảng; số lượng mặt hàng không đổi. & Đúng như mong đợi & \textbf{PASS}\\ \hline
\textbf{UC-TC-03} & SC-03 & Đang ở màn chỉnh sửa, đã sửa $\geq 1$ giá trị (dirty=true) & (1) Bấm ``Quay lại''. (2) Hộp thoại ``Thoát mà không lưu'' $\to$ chọn ``OK''. & Rời màn hình về danh sách; các thay đổi chưa lưu bị hủy (CSDL không đổi). & Đúng như mong đợi & \textbf{PASS}\\ \hline
\textbf{UC-TC-04} & SC-04 & Đang ở màn chỉnh sửa, đã sửa $\geq 1$ giá trị (dirty=true) & (1) Bấm ``Quay lại''. (2) Hộp thoại ``Thoát mà không lưu'' $\to$ chọn ``Cancel''. & Đóng hộp thoại; ở lại màn chỉnh sửa; mọi giá trị đã nhập được giữ nguyên. & Đúng như mong đợi & \textbf{PASS}\\ \hline
\textbf{UC-TC-05} & SC-05 & Mở màn chỉnh sửa, KHÔNG sửa gì (dirty=false) & (1) Bấm ``Quay lại''. & KHÔNG hiện hộp thoại xác nhận; rời màn hình ngay về danh sách. & Đúng như mong đợi & \textbf{PASS}\\ \hline
\textbf{UC-TC-06} & SC-06 & Yêu cầu 2002 ở trạng thái \texttt{processing} & (1) Mở màn chỉnh sửa 2002. & Thông báo ``Không thể chỉnh sửa'' (``...đang được xử lý nên không thể chỉnh sửa. Bạn chỉ có thể xem thông tin.''); nút ``Thêm mặt hàng'' ẩn; nút ``Lưu'' vô hiệu; ô SL \& ngày chỉ đọc; không có nút Xóa. & Đúng như mong đợi & \textbf{PASS}\\ \hline
\textbf{UC-TC-07} & SC-07 & Yêu cầu 3003 ở trạng thái \texttt{completed} & (1) Mở màn chỉnh sửa 3003. & Thông báo ``...đã hoàn tất nên không thể chỉnh sửa. Bạn chỉ có thể xem thông tin.''; giao diện chỉ đọc giống UC-TC-06. & Đúng như mong đợi & \textbf{PASS}\\ \hline
\textbf{UC-TC-08} & SC-08 & Màn hình ở chế độ không sửa được (editable=false) & (1) Kích hoạt hành động ``Lưu thay đổi'' (kiểm thử phòng vệ tầng controller). & Cảnh báo ``Yêu cầu [mã] không còn ở trạng thái PENDING, không thể lưu thay đổi.''; không ghi CSDL. & Đúng như mong đợi & \textbf{PASS}\\ \hline
\textbf{UC-TC-09} & SC-09 & Yêu cầu pending, ô Số lượng đang sửa được & (1) Kích ô Số lượng, nhập \texttt{-5}. (2) Nhấn Enter / rời ô. & Giá trị tự động chuyển về 0 (kẹp biên); không có giá trị âm nào được ghi nhận. & Đúng như mong đợi & \textbf{PASS}\\ \hline
\textbf{UC-TC-10} & SC-10 & Yêu cầu pending, ô Số lượng đang có giá trị 100 & (1) Nhập \texttt{abc} vào ô Số lượng. (2) Rời ô. & Ô tự khôi phục về giá trị số cũ (100); không thay đổi dữ liệu. & Đúng như mong đợi & \textbf{PASS}\\ \hline
\textbf{UC-TC-11} & SC-11 & Yêu cầu pending, đang mở hộp thoại ``Thêm mặt hàng'' & (1) Chọn 1 sản phẩm. (2) Gõ số lượng \texttt{0} (hoặc số âm \texttt{-3}). (3) Bấm ``Xác nhận''. & Hiện lỗi ``Số lượng phải là số nguyên dương.''; mặt hàng KHÔNG được thêm; hộp thoại vẫn mở. & Đúng như mong đợi & \textbf{PASS}\\ \hline
\textbf{UC-TC-12} & SC-12 & Yêu cầu pending, đang mở hộp thoại ``Thêm mặt hàng'' & (1) Chọn 1 sản phẩm. (2) Nhập SL hợp lệ (vd 10). (3) Xóa trống ô Ngày nhận. (4) Bấm ``Xác nhận''. & Hiện lỗi ``Vui lòng chọn ngày nhận mong muốn.''; mặt hàng KHÔNG được thêm. & Đúng như mong đợi & \textbf{PASS}\\ \hline
\textbf{UC-TC-13} & SC-13 & Không tồn tại yêu cầu mã 9999 trong CSDL & (1) Yêu cầu hệ thống mở màn chỉnh sửa mã 9999. & Cảnh báo ``Không thể tải yêu cầu'' kèm ``Không tìm thấy yêu cầu 9999''; không vào được màn chỉnh sửa với dữ liệu. & Đúng như mong đợi & \textbf{PASS}\\ \hline
\end{longtable}
}
\end{landscape}

\subsection{Tổng hợp kết quả kiểm thử Use Case}
\begin{table}[H]
\centering
\small
\begin{tabularx}{\textwidth}{|>{\bfseries}p{5.2cm}|Y|}
\hline
\rowcolor{gray!20}\normalfont\textbf{Chỉ số} & \textbf{Giá trị}\\ \hline
Tổng số kịch bản (scenarios) & 13 (1 chính + 4 thay thế + 8 ngoại lệ)\\ \hline
Tổng số test case use case & 13\\ \hline
Số test case PASS & \textbf{13 / 13}\\ \hline
Số test case FAIL & 0\\ \hline
Độ phủ luồng & Phủ 100\% các luồng đã đặc tả (Main Flow + 4 Alternative Flow + 5 Exception Flow)\\ \hline
\end{tabularx}
\end{table}

% =====================================================================
\section{Kết luận}
\begin{itemize}[leftmargin=*]
  \item \textbf{Số lượng Test Case đơn vị:} 9 (5 Hộp đen, 4 Hộp trắng). Đã chạy thực tế bằng JUnit 5: \textbf{9/9 PASS} (\texttt{Tests run: 9, Failures: 0, Errors: 0}).
  \item \textbf{Độ bao phủ nhánh (C1 coverage):} đạt \textbf{100\%} đối với phương thức \texttt{updateRequest}.
  \item \textbf{Số lượng Test Case Use Case:} 13 (1 kịch bản chính, 4 kịch bản thay thế, 8 kịch bản ngoại lệ). Kết quả: \textbf{13/13 PASS}.
  \item Hệ thống xử lý nghiệp vụ chỉnh sửa yêu cầu nhập hàng hoạt động ổn định, bắt đúng các ràng buộc nghiệp vụ (chỉ cho sửa khi \texttt{pending}), kiểm soát chặt dữ liệu nhập (số lượng, ngày nhận) và đồng bộ chính xác dữ liệu giữa UI, Service và CSDL.
\end{itemize}

\end{document}
